\chapter{Laparoscopy and Dye Test}

\section*{Overview}
Laparoscopy and dye test is a diagnostic fertility procedure in women that examines the uterus and fallopian tubes. A colored dye is passed through the uterus into the fallopian tubes while the organs are viewed with a laparoscope. It checks whether the tubes are open and whether sperm and egg can travel freely.

\section*{Alternative Names}
Dye test, chromoperturbation, hydrotubation with laparoscopy, tubal patency test with laparoscopy

\section*{Principle}
A thin telescope is inserted through a small incision in the navel to view the pelvis. Dye is then injected through the cervix into the uterus and fallopian tubes. If the dye flows out of the tube ends, the tubes are patent (open).

\section*{History}
Injection of dye to test tubal patency was described in the mid-1900s. Combining the test with laparoscopy allowed direct visualization and improved accuracy. It remains a standard part of the fertility workup.

\section*{Why It Is Done}
It is ordered as part of an infertility evaluation to determine whether the fallopian tubes are blocked. It is also used to assess the uterus, ovaries, and pelvic cavity. It may be recommended after pelvic infection, surgery, or endometriosis.

\section*{Is This a Primary or Secondary Test or Procedure}
It is a primary test for tubal patency in women who have not conceived. It is secondary when used to confirm findings seen on ultrasound or hysterosalpingography.

\section*{How It Is Performed}
The procedure is done under general anesthesia. After inflating the abdomen with gas, a laparoscope is inserted through the navel, and dye is injected through the cervix. The surgeon observes whether dye passes through each tube and also inspects the ovaries and pelvis.

\section*{Preparation}
The test is timed for the first half of the menstrual cycle, before ovulation. Fasting for several hours is required because general anesthesia is used. Pregnancy is excluded beforehand.

\section*{Contraindications and Precautions}
\begin{itemize}
  \item Confirmed or suspected pregnancy
  \item Active pelvic infection
  \item Allergy to the dye used
  \item Heavy bleeding during the current menstrual period
\end{itemize}

\section*{Risks}
Bleeding, infection, and injury to pelvic organs are possible. Shoulder and abdominal pain may occur from the gas used to inflate the abdomen. Minor adhesions may be treated at the same time.

\section*{Aftercare and Recovery}
Most women go home the same day after recovery from anesthesia. Mild cramping, bloating, and a little vaginal spotting are common. Strenuous activity is usually avoided for a few days.

\section*{Outcome and Follow-up}
If dye flows freely from both tubes, the tubes are patent and no tubal cause for infertility is found. If dye does not pass, the tube is blocked, and the level and site of blockage guide treatment.

\section*{Expected Results}
Free flow of dye through both fallopian tubes, confirming patency. The uterus, ovaries, and pelvic lining appear normal.

\section*{Follow-up}
Results are discussed with a fertility specialist who plans the next steps, such as timed intercourse or assisted reproduction. If blockages are found, surgery or in vitro fertilization may be considered. Repeat testing is only done if symptoms suggest new problems.

\section*{When to Seek Medical Care}
Follow the procedure team's specific instructions. Seek urgent medical assessment for severe or worsening pain, uncontrolled bleeding, fever or signs of infection, trouble breathing, chest pain, fainting, new weakness or confusion, inability to urinate, or any rapidly worsening symptom. The appropriate warning signs depend on the procedure and the patient's medical condition.

\section*{References}
\begin{enumerate}
  \item MedlinePlus. Laparoscopy and Dye Test. U.S. National Library of Medicine; 2025.
  \item Current Medical Diagnosis and Treatment. McGraw-Hill; 2025.
\end{enumerate}

\clearpage
